\chapter*{Lời Mở Đầu}
\addcontentsline{toc}{chapter}{Lời Mở Đầu}
\markboth{Lời Mở Đầu}{Lời Mở Đầu}

\begin{truyen}{Lời Mở Đầu}{Opening Words}
\chuhoa{K}{hông ai phát cho em một bản nháp trước khi sống.}
\textit{(No one hands you a draft before you start living.)}

Rời ghế nhà trường là một ngưỡng cửa mà hầu hết ai cũng bước qua chưa chuẩn. Công việc đầu tiên, tiền đầu tiên, bị từ chối, bị thay thế, sống xa nhà, bạn bè ra rẽ — tất cả đều là những tình huống chưa có trong sách giáo khoa.
\textit{(Leaving school is a threshold that most people cross unprepared. First job, first paycheck, rejection, replacement, living away from home, friends parting ways --- none of these are in textbooks.)}

Quyển chín không hứa hẹn chuẩn bị cho em đầy đủ. Không ai có thể làm được điều đó. Nhưng nó có thể giúp em bớt ngây thơ, bớt bỡ ngỡ, và biết rằng những gì em đang trải qua là bình thường.
\textit{(Volume nine does not promise to fully prepare you. No one can do that. But it can help you be less naive, less bewildered, and know that what you are experiencing is normal.)}

Tỉnh táo cũng là một kiểu tự cứu mình.
\textit{(Being clear-eyed is also a way of saving yourself.)}
\end{truyen}

\begin{baihoc}
Ra đời không có bản nháp --- nhưng có thể học từ những người đã bước trước.
\textit{(Life has no draft --- but we can learn from those who walked ahead.)}
\end{baihoc}

\begin{ghinhoanh}
\textbf{Short English to remember:}

Life begins without a rehearsal.

What you feel stepping out is normal.
\end{ghinhoanh}

\begin{tuhocnhanh}
1. Điều gì em lo lắng nhất khi nghĩ đến ngày rời trường? \textit{What worries you most when you think about the day you leave school?}

2. Em cảm thấy mình đang được chuẩn bị cho cuộc đời chưa? \textit{Do you feel prepared for life so far?}

3. Bài học nào trong quyển này em muốn ghi nhớ trước khi bước vào đời? \textit{Which lesson in this volume do you want to remember before stepping into the world?}
\end{tuhocnhanh}

\ngancach
